%% LyX 2.5.1 created this file.  For more info, see https://www.lyx.org/.
%% Do not edit unless you really know what you are doing.
\documentclass[american]{article}
\usepackage[T1]{fontenc}
\usepackage[utf8]{inputenc}
\usepackage{parskip}
\usepackage{babel}
\begin{document}
\title{A Brief Derivation of the Corrected Welch-Satterthwaite Equation}
\author{Matthias von Davier}
\date{July 2026}
\maketitle

\section{Introduction}

Assume $k$ independent variance components $S^{2}_{k}$ with $\frac{\nu_{k}S^{2}_{k}}{\sigma^{2}_{k}}\sim\chi^{2}_{\nu_{k}}$
for $1,\dots,K$. Typically, one may assume $S^{2}_{k}=\nu^{-1}_{k}\sum^{\nu_{k}}_{i=0}\left(x_{ik}-\overline{X_{k}}\right)^{2}$with
$\overline{X_{k}}=\frac{1}{\nu_{k}+1}\sum^{\nu_{k}}_{i=0}x_{ik}$

and $X_{ik}\sim N\left(\nu_{k},\sigma_{k}\right)$. Then define
\[
S^{2}=\sum^{K}_{k=1}w_{k}S^{2}_{k}
\]
with weights $w_{k}>0$. Important cases are variance estimates based
on combinations of imputation and sampling variance (REF) as well
as variance estimates based on resampling schemes such as the jackknife
(REF) and balanced repeated replication (FAY).

Satterthwaite (1941) conjectures that 
\[
\frac{\nu^{*}S^{2}}{\sigma^{2}}\sim\chi^{2}_{\nu^{*}}
\]
with unknown degrees of freedom $\nu^{*}$ and with variance $\sigma^{2}=E\left(S^{2}\right).$
Welch (1947) discusses a special case with $K=2$ as an improvement
over the $t-$test (Student, 1921).

\section{Derivation}

The assuming of independence of the variance components $S^{2}_{K}$
allows equating 
\[
V\left(S^{2}\right)=\sum^{K}_{k=1}w^{2}_{i}V\left(S^{2}_{k}\right).
\]

For $\chi^{2}$ random variables we have $V\left[\frac{\nu_{k}S^{2}_{k}}{\sigma^{2}_{k}}\right]=2\nu_{k}=\left[\frac{\nu_{k}}{\sigma^{2}_{k}}\right]^{2}V\left(S^{2}_{k}\right)$
and hence $V\left(S^{2}_{k}\right)=2\frac{\left[\sigma^{2}_{k}\right]^{2}}{\nu_{k}}$
and the same for $V\left(S^{2}\right)=2\frac{\left[\sigma^{2}\right]^{2}}{\nu^{*}}$.
Then we have 
\[
2\frac{\left[\sigma^{2}\right]^{2}}{\nu^{*}}=\sum^{K}_{k=1}2w^{2}_{k}\frac{\left[\sigma^{2}_{k}\right]^{2}}{\nu_{k}}.
\]
A well known property of the variance is $V\left(U\right)=E\left(U^{2}\right)-\left[E\left(U\right)\right]^{2}$,
which implies $E\left(U^{2}\right)=V\left(U\right)+\left[E\left(U\right)\right]^{2}$.
For $U\sim\chi^{2}_{\nu}$ , this yields $E\left(U^{2}\right)=2\nu+\nu^{2}=\nu\left(2+\nu\right).$

Hence 
\[
E\left(\left[\frac{\nu_{k}S^{2}_{k}}{\sigma^{2}_{k}}\right]^{2}\right)=\left[\frac{\nu_{k}}{\sigma^{2}_{k}}\right]^{2}E\left(\left[S^{2}_{k}\right]^{2}\right)=\nu_{k}\left[2+\nu_{k}\right]
\]
and further
\[
E\left(\left[S^{2}_{k}\right]^{2}\right)\frac{\nu_{k}}{2+\nu_{k}}=\left[\sigma^{2}_{k}\right]^{2}
\]
and the same for $S^{2},$i.e., $E\left(\left[S^{2}\right]^{2}\right)\frac{\nu^{*}}{2+\nu^{*}}=\left[\sigma^{2}\right]^{2}.$

\section{Plug-In Estimator}

Plugging in $\left[S^{2}_{k}\right]^{2}$for $E\left(\left[S^{2}_{k}\right]^{2}\right)$
and $\left[S^{2}\right]^{2}$ for $E\left(\left[S^{2}\right]^{2}\right)$
yields
\[
\frac{\left[S^{2}\right]^{2}\frac{\nu^{*}}{2+\nu^{*}}}{\nu^{*}}=\sum^{K}_{k=1}w^{2}_{k}\frac{\left[S^{2}_{k}\right]^{2}\frac{\nu_{k}}{2+\nu_{k}}}{\nu_{k}}
\]
and solving for $\nu^{*}$yields
\[
\nu^{*}_{vD}=\frac{\left[\sum^{K}_{k=1}w_{k}S^{2}_{k}\right]^{2}}{\sum^{K}_{k=1}w^{2}_{k}\frac{\left[S^{2}_{k}\right]^{2}}{\nu_{k}+2}}-2
\]
 which is the estimator proposed by von Davier (2026), also referred
to as corrected Welch-Satterthwaite equation.

\section{Original Large Sample Estimator}

Satterthwaite directly plugs in $\left[S^{2}\right]^{2}$ for $\left[\sigma^{2}\right]^{2}$and
$\left[S^{2}_{k}\right]^{2}$ for $\left[\sigma^{2}_{k}\right]^{2}$
thus effectively omitting the factor $\frac{\nu}{\nu+2}$ in the process.
The Welch-Satterthwaite equation is therefore given by
\[
\nu^{*}_{WS}=\frac{\left[\sum^{K}_{k=1}w_{k}S^{2}_{k}\right]^{2}}{\sum^{K}_{k=1}w^{2}_{k}\frac{\left[S^{2}_{k}\right]^{2}}{\nu_{k}}}
\]
and it is known to produce heavy downward bias for small $\nu_{k}$
(e.g., Johnson \& Rust, 1992; REFS). It is unclear whether this omission
was intentional or an oversight. However, Welch (1947) also omits
the factor $\frac{\nu}{\nu+2}$. The only instance known to the author
where this factor is used is a simulation study by Boardman (1974)
where one of the candidates for estimating $\nu^{*}$ is 
\[
\nu^{*}_{Boardman}=\frac{\left[\sum^{K}_{k=1}w_{k}S^{2}_{k}\right]^{2}}{\sum^{K}_{k=1}w^{2}_{k}\frac{\left[S^{2}_{k}\right]^{2}}{\nu_{k}+2}}
\]
which he calls the \emph{unbiased Satterthwaite}. Note that $\nu^{*}_{Boardman}=\nu^{*}_{vD}+2$.
von Davier (2026) showed that his proposal has a small upward bias
for very small $\nu_{k}(=1,2,3)$, that quickly vanishes with growing
$\nu_{k}.$ Boardman's (1974) is therefore affected by a larger upward
bias as it is larger by 2.

\section{Characteristics of the Corrected Degrees of Freedom Estimator}

\subsection{Minimum $\nu^{*}$}

Note that the minimum $\nu^{*}$ is attained when $S^{2}_{k}=0$ for
all but one $S^{2}_{l}>0$ with $\forall k:\nu_{l}\le\nu_{k}$. In
that case we have
\[
\left[\sum^{K}_{k=1}w_{i}S^{2}_{k}\right]^{2}=w^{2}_{l}\left[S^{2}_{l}\right]
\]
and 
\[
\sum^{K}_{k=1}w^{2}_{k}\frac{\left[S^{2}_{k}\right]^{2}}{\nu_{k}+2}=w^{2}_{k'}\frac{\left[S^{2}_{l}\right]^{2}}{\nu_{l}+2}
\]
 and hence 
\[
\nu^{*}_{vD}=\frac{\left[\sum^{K}_{k=1}w_{i}S^{2}_{k}\right]^{2}}{\sum^{K}_{k=1}w^{2}_{k}\frac{\left[S^{2}_{k}\right]^{2}}{\nu_{k}+2}}-2=\frac{w^{2}_{k'}\left[S^{2}_{l}\right]^{2}}{w^{2}_{k'}\frac{\left[S^{2}_{k'}\right]^{2}}{\nu_{l}+2}}-2=\nu_{l}>0.
\]

This implies that the estimator is non-negative. 

\subsection{Maximum $\nu^{*}$}

The maximum is attained if $w_{k}S^{2}_{k}=w_{l}S^{2}_{l}=C$ for
all $k,l=1,\dots,K.$ In that case, we have
\[
\nu^{*}_{vD}=\frac{\left[\sum^{K}_{k=1}C\right]^{2}}{\sum^{K}_{k=1}\frac{C^{2}}{\nu_{k}+2}}-2=\frac{K^{2}}{\sum^{K}_{k=1}\frac{1}{\nu_{k}+2}}-2=K\times H\left(\nu_{1}+2,\dots.,\nu_{K}+2\right)-2.
\]
This expression is a function of the harmonic mean of the component
degrees of freedom (plus 2) and the number of components $K$.

\subsection{Known and Equal Variance Components $S^{2}_{k}\equiv\sigma^{2}_{E}$
for all $k$}

If the variance components are all equal and known, $S^{2}_{k}\equiv\sigma^{2}_{E}$,
we obtain
\[
S^{2}=\sigma^{2}=K\sigma^{2}_{E}.
\]
In this case, no estimator of variance components is used. Then, equality
of all components provides 
\[
\nu^{*}=\frac{\left(\sum^{K}_{k=1}w_{k}\right)^{2}}{\sum^{K}_{k=1}w^{2}_{k}}=N_{eff}
\]
which is the effective sample size estimator proposed by Kish (1965).

\section*{References}

Boardman, T. J. (1974). Confidence Intervals for Variance Components
-- A Comparative Monte Carlo Study. Biometrics, 30(2), 251--262.
doi:10.2307/2529647

Johnson, E. G., \& Rust, K. F. (1992). Population Inferences and Variance
Estimation for NAEP Data. Journal of Educational Statistics, 17(2),
175--190. doi:10.2307/1165168

Kish, L. (1965). Survey sampling. John Wiley \& Sons.

Kish, L. (1992). Weighting for unequal Pi. Journal of Official Statistics,
8(2), 183--200.

Lipsitz, S., Parzen, M., \& Zhao, L. P. (2002). A Degrees-Of-Freedom
approximation in multiple imputation. Journal of Statistical Computation
and Simulation, 72(4), 309--318. doi:10.1080/00949650212848

Rubin, D. B., \& Schenker, N. (1986). Multiple Imputation for Interval
Estimation From Simple Random Samples With Ignorable Nonresponse.
Journal of the American Statistical Association, 81(394), 366--374.
doi:10.2307/2289225

Satterthwaite, F. E. (1941). Synthesis of Variance. Psychometrika,
6(5), 309-316. doi:10.1007/BF02288586

Satterthwaite, F. E. (1946). An Approximate Distribution of Estimates
of Variance Components. Biometrics Bulletin, 2(6), 110--114. doi:10.2307/3002019

von Davier, M. (2025). An Improved Satterthwaite (1941, 1946) Effective
df Approximation. Journal of Educational and Behavioral Statistics,
0(0). doi:10.3102/10769986241309329

Welch, B. L. (1947). The generalization of ‘Student’s’ problem when
several different population variances are involved. Biometrika, 34(1/2),
28-35. doi:10.2307/2332510
\end{document}
